\section{实验结果与分析}
\label{sec:results-discussion}

本章按实验设置、主对比、消融、显著性和开销分析的顺序报告结果。所有结论均基于既有正式结果包，不重跑优化实验；单一标量适应度不能完全代表强约束低空路径质量，因此下文结合可行率、平均排名、工程轨迹复核标记、代表性路径和评价开销进行解释。

\subsection{整体实验设置}
\label{sec:experimental-setup}

实验在同一台本地工作站上完成：Intel(R) Core(TM) i5-14600KF \SI{3.50}{GHz}，\SI{32}{GB} 内存，64 位系统，MATLAB R2022b。所有正式运行均为单进程串行；由于 wall-clock time 受 MATLAB 会话状态、repair 次数和批处理负载影响，效率主要由函数评价次数、CVF 触发规模和同批次评价比值刻画。

\begin{table}[tbp]
\centering
\caption{三维低空路径规划测试场景设置}
\label{tab:scene-setup}
\small
\begin{tabular}{@{}lcccc@{}}
\toprule
场景 & 障碍物数 & 禁飞区数 & 风险热点数 & 高度范围 \\
\midrule
Scene 1 & 6 & 3 & 3 & $[8,40]$ \\
Scene 2 & 10 & 2 & 3 & $[8,32]$ \\
Scene 3 & 17 & 3 & 5 & $[8,26]$ \\
\bottomrule
\end{tabular}
\end{table}

三维规划空间水平方向均为 $[0,100]\times[0,100]$，起点和终点分别为 $\mathbf{s}=[5,5,12]^\mathsf{T}$ 和 $\mathbf{g}=[95,95,18]^\mathsf{T}$。三个场景均含建筑物、柱状禁飞区和风险热点，约束密度、高度上界和风险分布逐步收紧；Scene 1/2 使用 $z_{\mathrm{ref}}=16$，Scene 3 使用 $z_{\mathrm{ref}}=14$。

\begin{table}[tbp]
\centering
\caption{公共路径编码与评价参数设置}
\label{tab:common-setup}
\small
\begin{tabularx}{\linewidth}{@{}l X@{}}
\toprule
类别 & 设置 \\
\midrule
路径编码 & $K=6$，维度 $18$ \\
B 样条 & 三次 B 样条，$M=220$ \\
搜索预算 & $N=30$，$T=300$，30 次独立运行 \\
目标权重 & $w_L=1.00$，$w_E=0.60$，$w_R=1.20$，$w_S=0.15$，$w_H=0.80$，$w_B=0.80$ \\
违反惩罚 & $\lambda_{\mathrm{obs}}=150$，$\lambda_{\mathrm{nfz}}=200$，$\lambda_{\mathrm{curv}}=20$，$\lambda_{\mathrm{alt}}=30$ \\
运动约束 & $\theta_{\max}=55^\circ$，默认高度 $[8,40]$ \\
\bottomrule
\end{tabularx}
\end{table}

基线选择遵循覆盖性和可复现性原则：AE 作为基础进化搜索骨架，PSO、GWO 和 HHO 代表经典群智能优化方法，ERIME 和 MSCSO 代表近期改进型元启发式方法。所有算法使用相同路径编码、评价函数、边界投影和 Deb 可行性优先准则；基线不调用 CVF 场、约束状态调度、稀疏可行性保持或 CVF-AE 专属初始化。消融实验与主对比保持相同场景、运行次数、种群规模和迭代次数。

主要指标包括最终最佳适应度、可行率、最终违反量、函数评价次数、平均排名和工程轨迹复核标记。工程复核规则由高空绕行、边界贴靠、过度绕行和控制点贴边四类固定阈值标记取逻辑或得到，仅作 sanity check，不作为安全认证指标或独立优化目标。由于不同算法的随机种子不视为严格配对样本，CVF-AE 与各基线采用双侧 Wilcoxon rank-sum / Mann--Whitney U 检验 \citep{Mann1947RankSum}，场景内多重比较采用 Holm 校正 \citep{Holm1979Sequential}，并报告面向最小化问题定义的 Cliff's delta \citep{Cliff1993Dominance}。


\begin{figure}[tbp]
\centering
\includegraphics[width=\linewidth]{selected7_boxplots.png}
\caption{主对比算法的最终适应度分布}
\label{fig:boxplots-strong}
\end{figure}
\FloatBarrier
\subsection{主对比结果}

\begin{table*}[tbp]
\centering
\caption{主对比实验结果}
\label{tab:main-comparison}
\scriptsize
\setlength{\tabcolsep}{3pt}
\resizebox{\textwidth}{!}{%
\begin{tabular}{lccccccccccc}
\toprule
Algorithm & S1 $J$ & S1 FR & S1 R & S2 $J$ & S2 FR & S2 R & S3 $J$ & S3 FR & S3 R & Avg. R & Eng. flag \\
\midrule
AE & 656.1972 $\pm$ 325.9113 & 0.60 & 7 & 1627.2324 $\pm$ 330.2622 & 0.00 & 7 & 1680.4468 $\pm$ 208.6819 & 0.00 & 7 & 7.00 & 0.22 \\
PSO & 423.3491 $\pm$ 216.2313 & 0.90 & 6 & 659.5587 $\pm$ 303.2289 & 0.43 & 6 & 1303.4640 $\pm$ 388.2979 & 0.03 & 6 & 6.00 & 0.49 \\
GWO & 335.7826 $\pm$ 32.9433 & 0.97 & 4 & 341.1693 $\pm$ 41.4608 & 1.00 & 4 & 673.2786 $\pm$ 376.6002 & 0.53 & 4 & 4.00 & 0.93 \\
HHO & 361.7454 $\pm$ 71.7330 & 0.97 & 5 & 357.9223 $\pm$ 149.3598 & 0.97 & 5 & 929.5468 $\pm$ 437.4162 & 0.23 & 5 & 5.00 & 0.78 \\
ERIME & 318.5440 $\pm$ 15.4114 & 1.00 & 3 & 320.8775 $\pm$ 24.5106 & 1.00 & 3 & 356.4928 $\pm$ 24.6010 & 1.00 & 2 & 2.67 & 0.79 \\
MSCSO & 307.8580 $\pm$ 14.8041 & 1.00 & 2 & 295.0138 $\pm$ 19.1301 & 1.00 & 1 & 449.9550 $\pm$ 198.8820 & 0.83 & 3 & 2.00 & 0.64 \\
CVF-AE & 298.5286 $\pm$ 2.6211 & 1.00 & 1 & 301.1934 $\pm$ 2.8253 & 1.00 & 2 & 350.7658 $\pm$ 23.3272 & 1.00 & 1 & 1.33 & 0.10 \\
\bottomrule
\end{tabular}%
}
\vspace{1mm}
\parbox{\textwidth}{\footnotesize 注：FR 为可行率，R 为主对比算法按平均适应度升序得到的场景排名，Avg. R 为三场景平均排名。Eng. flag 为三场景共 90 次运行中被任一工程轨迹复核规则标记的比例；该固定阈值复核仅作 sanity check，不作为安全认证指标。}
\end{table*}


表 \ref{tab:main-comparison} 给出了主对比结果。CVF-AE 在三个场景中的可行率均为 $1.00$，平均场景排名为 $1.33$，并取得所选代表性基线中最佳平均排名；其工程轨迹复核标记比例为 $0.10$，优势主要体现为可行性、平均排名和轨迹形态诊断之间的折中。



图 \ref{fig:boxplots-strong} 展示逐次运行分布。CVF-AE 在 Scene 1 和 Scene 3 中箱体较窄且均值最低；Scene 2 中 MSCSO 的标量适应度更低，但其工程轨迹复核标记比例较高，因此该场景不支持 CVF-AE 标量最优结论，而支持其更保守的低空走廊折中。Scene 3 中，CVF-AE 的平均适应度为 $350.7658$，略低于 ERIME 的 $356.4928$ 和 MSCSO 的 $449.9550$，并保持 $1.00$ 可行率；ERIME 虽保持可行率，但工程复核标记比例较高。

\begin{figure}[tbp]
\centering
\begin{subfigure}{0.96\linewidth}
    \centering
    \includegraphics[width=\linewidth]{selected_scene1_convergence_mean_std.png}
    \caption{Scene 1}
\end{subfigure}

\begin{subfigure}{0.96\linewidth}
    \centering
    \includegraphics[width=\linewidth]{selected_scene2_convergence_mean_std.png}
    \caption{Scene 2}
\end{subfigure}

\begin{subfigure}{0.96\linewidth}
    \centering
    \includegraphics[width=\linewidth]{selected_scene3_convergence_mean_std.png}
    \caption{Scene 3}
\end{subfigure}
\caption{三类测试场景的平均收敛曲线}
\label{fig:convergence}
\end{figure}

图 \ref{fig:convergence} 给出平均收敛行为。CVF-AE 在 Scene 1 和 Scene 3 中较快进入稳定区间；Scene 2 中 MSCSO 最终均值更低，对应主表中的标量适应度优势，但仍需结合轨迹形态解释。

\begin{figure*}[tbp]
\centering
\begin{subfigure}{0.32\linewidth}
    \centering
    \includegraphics[width=\linewidth]{selected_scene1_representative_top.png}
    \caption{Scene 1}
\end{subfigure}
\hfill
\begin{subfigure}{0.32\linewidth}
    \centering
    \includegraphics[width=\linewidth]{selected_scene2_representative_top.png}
    \caption{Scene 2}
\end{subfigure}
\hfill
\begin{subfigure}{0.32\linewidth}
    \centering
    \includegraphics[width=\linewidth]{selected_scene3_representative_top.png}
    \caption{Scene 3}
\end{subfigure}
\caption{主对比算法的路径俯视对比}
\label{fig:selected-paths-top}
\end{figure*}

\begin{figure*}[tbp]
\centering
\begin{subfigure}{0.32\linewidth}
    \centering
    \includegraphics[width=\linewidth]{selected_scene1_representative_3d.png}
    \caption{Scene 1}
\end{subfigure}
\hfill
\begin{subfigure}{0.32\linewidth}
    \centering
    \includegraphics[width=\linewidth]{selected_scene2_representative_3d.png}
    \caption{Scene 2}
\end{subfigure}
\hfill
\begin{subfigure}{0.32\linewidth}
    \centering
    \includegraphics[width=\linewidth]{selected_scene3_representative_3d.png}
    \caption{Scene 3}
\end{subfigure}
\caption{主对比算法的路径三维对比}
\label{fig:selected-paths-3d}
\end{figure*}

图 \ref{fig:selected-paths-top} 和图 \ref{fig:selected-paths-3d} 解释上述数值差异。代表性路径仅用于可视化展示，不参与优化、统计检验或表格排名；Scene 2 俯视图显示 CVF-AE 避开建筑足迹并沿走廊推进，Scene 3 中部分基线在低高度和密集障碍下出现边界贴靠或大幅绕行，说明强约束低空路径规划不能只按最终适应度排序。

% \FloatBarrier
\subsection{消融分析}

为压缩表格，消融版本记为 V0--V5：V0 为完整 CVF-AE，V1 去除 CVF 场，V2 去除状态自适应 CVF，V3 去除稀疏可行性保持，V4 去除约束感知初始化，V5 为 Base-AE。

\begin{table*}[tbp]
\centering
\caption{CVF-AE 关键模块消融实验结果}
\label{tab:ablation}
\setlength{\tabcolsep}{3pt}
\tiny
\resizebox{\textwidth}{!}{%
\begin{tabular}{clcccccc}
\toprule
Scene & Variant & $J$ & FR & R & Avg. R & nEvals & CVF trig./succ. \\
\midrule
1 & CVF-AE & 297.8318 $\pm$ 3.1119 & 1.00 & 1 & 1.00 & 9566.4 & 21.1/19.2 \\
1 & CVF-AE-w/o-CVF & 306.1985 $\pm$ 3.4917 & 1.00 & 3 & 3.00 & 9471.8 & 0.0/0.0 \\
1 & CVF-AE-w/o-StateAdaptiveCVF & 298.3037 $\pm$ 3.2161 & 1.00 & 2 & 2.33 & 9566.0 & 20.6/19.7 \\
1 & CVF-AE-w/o-SparsePreservation & 306.3793 $\pm$ 4.9973 & 1.00 & 4 & 3.67 & 9140.1 & 110.1/90.7 \\
1 & CVF-AE-w/o-Init & 316.9964 $\pm$ 11.8481 & 1.00 & 5 & 5.00 & 9587.6 & 140.2/93.1 \\
1 & Base-AE & 654.0475 $\pm$ 328.7142 & 0.63 & 6 & 6.00 & 9030.0 & 0.0/0.0 \\
2 & CVF-AE & 300.6612 $\pm$ 3.2775 & 1.00 & 1 & 1.00 & 9551.6 & 10.6/7.9 \\
2 & CVF-AE-w/o-CVF & 303.7697 $\pm$ 4.9565 & 1.00 & 4 & 3.00 & 9432.5 & 0.0/0.0 \\
2 & CVF-AE-w/o-StateAdaptiveCVF & 301.8164 $\pm$ 2.9332 & 1.00 & 2 & 2.33 & 9560.9 & 20.4/14.6 \\
2 & CVF-AE-w/o-SparsePreservation & 303.3114 $\pm$ 6.0718 & 1.00 & 3 & 3.67 & 9107.8 & 77.8/51.5 \\
2 & CVF-AE-w/o-Init & 313.2191 $\pm$ 18.3412 & 1.00 & 5 & 5.00 & 9592.4 & 153.7/107.4 \\
2 & Base-AE & 1597.4829 $\pm$ 342.4791 & 0.00 & 6 & 6.00 & 9030.0 & 0.0/0.0 \\
3 & CVF-AE & 348.5856 $\pm$ 21.8153 & 1.00 & 1 & 1.00 & 9367.1 & 86.8/59.5 \\
3 & CVF-AE-w/o-CVF & 355.8347 $\pm$ 21.3086 & 1.00 & 2 & 3.00 & 9258.5 & 0.0/0.0 \\
3 & CVF-AE-w/o-StateAdaptiveCVF & 361.3123 $\pm$ 35.0403 & 1.00 & 3 & 2.33 & 9321.8 & 72.5/46.3 \\
3 & CVF-AE-w/o-SparsePreservation & 372.2255 $\pm$ 19.8216 & 1.00 & 4 & 3.67 & 9118.7 & 88.7/56.6 \\
3 & CVF-AE-w/o-Init & 556.7515 $\pm$ 348.8264 & 0.70 & 5 & 5.00 & 9802.7 & 417.9/336.5 \\
3 & Base-AE & 1643.7320 $\pm$ 199.4027 & 0.00 & 6 & 6.00 & 9030.0 & 0.0/0.0 \\
\bottomrule
\end{tabular}%
}
\vspace{1mm}
\parbox{\textwidth}{\footnotesize 注：FR 为可行率，R 为场景内排名，Avg. R 为跨场景平均排名。CVF trig./succ. 表示平均触发的 CVF 候选数 / 平均被局部接收的 CVF 候选数。}
\end{table*}


表 \ref{tab:ablation} 展示正式消融结果。完整 CVF-AE 的消融平均排名为 $1.00$，并在所有场景保持 $1.00$ 可行率；相对 Base-AE，它在 Scene 2/3 中将可行率从 $0$ 提升到 $1.00$，说明约束感知初始化、CVF 候选和稀疏可行性保持共同支撑可行路径形成。

\begin{figure}[tbp]
\centering
\includegraphics[width=0.78\linewidth]{ablation_scene2_representative_top.png}
\caption{Scene 2 关键模块消融版本的路径俯视对比}
\label{fig:ablation-scene2-top}
\end{figure}

\begin{figure}[tbp]
\centering
\includegraphics[width=0.78\linewidth]{ablation_scene3_representative_3d.png}
\caption{Scene 3 关键模块消融版本的路径三维对比}
\label{fig:ablation-paths}
\label{fig:ablation-scene3-3d}
\end{figure}

图 \ref{fig:ablation-scene2-top} 和图 \ref{fig:ablation-scene3-3d} 展示消融路径形态。Scene 2 俯视图显示完整方法更稳定地沿可行走廊推进；Scene 3 三维图说明，在更低高度上界和更密集约束下，去除初始化或 CVF 相关机制后更容易出现低质量绕行。去除 CVF 场后算法仍可保持可行，但适应度上升：Scene 1 从 $297.8318$ 到 $306.1985$，Scene 2 从 $300.6612$ 到 $303.7697$，Scene 3 从 $348.5856$ 到 $355.8347$。去除约束感知初始化在 Scene 3 影响最大，可行率从 $1.00$ 降至 $0.70$，平均适应度升至 $556.7515$。

% \FloatBarrier
\subsection{显著性分析}

\begin{table}[tbp]
\centering
\caption{CVF-AE 与基线算法的逐场景显著性比较}
\label{tab:significance}
\scriptsize
\setlength{\tabcolsep}{2.8pt}
\begin{tabular}{@{}llccc@{}}
\toprule
Baseline & Scene & $p_{\mathrm{Holm}}$ & $\delta$ & R \\
\midrule
AE & 1 & 3.0e-10 & +1.00 & + \\
AE & 2 & 3.0e-10 & +1.00 & + \\
AE & 3 & 3.0e-10 & +1.00 & + \\
PSO & 1 & 2.0e-07 & +0.81 & + \\
PSO & 2 & 5.5e-04 & +0.59 & + \\
PSO & 3 & 3.0e-10 & +1.00 & + \\
GWO & 1 & 3.0e-10 & +1.00 & + \\
GWO & 2 & 0.001 & +0.55 & + \\
GWO & 3 & 2.1e-05 & +0.70 & + \\
WOA & 1 & 3.0e-10 & +1.00 & + \\
WOA & 2 & 3.0e-10 & +1.00 & + \\
WOA & 3 & 3.0e-10 & +1.00 & + \\
HHO & 1 & 1.0e-09 & +0.95 & + \\
HHO & 2 & 0.001 & +0.55 & + \\
HHO & 3 & 5.9e-08 & +0.87 & + \\
DBO & 1 & 3.0e-10 & +1.00 & + \\
DBO & 2 & 1.5e-06 & +0.78 & + \\
DBO & 3 & 5.1e-05 & +0.66 & + \\
CPO & 1 & 0.773 & +0.04 & = \\
CPO & 2 & 3.0e-10 & -1.00 & - \\
CPO & 3 & 0.934 & +0.14 & = \\
GDSAO & 1 & 5.5e-10 & +0.97 & + \\
GDSAO & 2 & 0.620 & -0.08 & = \\
GDSAO & 3 & 0.934 & +0.15 & = \\
ERIME & 1 & 4.0e-10 & +0.98 & + \\
ERIME & 2 & 0.595 & +0.16 & = \\
ERIME & 3 & 0.097 & +0.34 & = \\
MSCSO & 1 & 4.8e-07 & +0.78 & + \\
MSCSO & 2 & 0.002 & -0.51 & - \\
MSCSO & 3 & 0.934 & +0.13 & = \\
\bottomrule
\end{tabular}
\end{table}


表 \ref{tab:significance} 给出了 CVF-AE 与 6 个基线的逐场景 Holm 校正 $p$ 值和 Cliff's delta。Holm 校正范围限定为每个场景内 6 次比较；成对比较中，CVF-AE 在 $18$ 次比较中得到 $13$ 胜、$4$ 平和 $1$ 负。唯一显著劣势出现在 Scene 2，对应 MSCSO 的标量适应度优势。表中 $R$ 为 Holm 校正后的胜、平、负关系，$\delta>0$ 表示 CVF-AE 分布更优。统计检验只作用于最终最佳适应度分布，不能替代工程轨迹复核，因此本文将结论限定为综合均衡优势，而非逐场景标量最优。

% \FloatBarrier
\subsection{开销分析}

\begin{table}[tbp]
\centering
\caption{CVF-AE 评价开销与稀疏触发统计}
\label{tab:runtime}
\tiny
\resizebox{\columnwidth}{!}{%
\begin{tabular}{@{}lcccccc@{}}
\toprule
Group & Sc. & nEvals & Trig. & Succ. & Rate & Ratio \\
\midrule
Main & 1 & 9569.5 & 25.0 & 22.4 & 0.28\% & -- \\
Main & 2 & 9552.7 & 12.4 & 9.5 & 0.14\% & -- \\
Main & 3 & 9362.5 & 83.8 & 56.4 & 0.93\% & -- \\
\midrule
Abl. & 1 & 9566.4 & 21.1 & 19.2 & 0.22\% & 1.010 \\
Abl. & 2 & 9551.6 & 10.6 & 7.9 & 0.11\% & 1.013 \\
Abl. & 3 & 9367.1 & 86.8 & 59.5 & 0.93\% & 1.012 \\
\bottomrule
\end{tabular}%
}
\vspace{1mm}
\parbox{\columnwidth}{\footnotesize 注：Trig. 和 Succ. 分别表示平均触发的 CVF 候选数和平均被接收的 CVF 候选数；Ratio 为完整 CVF-AE 相对于同批次 w/o-CVF 消融版本的评价次数比。}
\end{table}


表 \ref{tab:runtime} 给出了评价开销摘要。CVF-AE 在 Scene 1--3 的平均评价次数分别为 $9569.5$、$9552.7$ 和 $9362.5$，与 $N=30$、$T=300$ 的基础规模接近。CVF 触发次数分别为 $25.0$、$12.4$ 和 $83.8$；更强约束场景触发更多，但触发率仍低于 $1\%$。同批次消融中，完整 CVF-AE 相对于去除 CVF 的评价次数比值约为 $1.010$、$1.013$ 和 $1.012$，说明额外评价开销约为 $1\%$ 量级。本文不以绝对 wall-clock time 作为主要效率证据。

参数稳健性检查仅作补充。更高触发配额可在 Scene 2/3 降低部分适应度均值，但会增加触发和评价规模；步长强度较优档位随场景变化；状态切换阈值未出现同时稳定占优的非默认设置。因此，本文保留默认参数，并将结论限定为默认配置的折中稳健性。

综合主对比、消融、显著性和开销分析，CVF-AE 的价值在于：在所选代表性基线中保持 $1.00$ 可行率并取得最佳平均排名；在 Scene 2/3 提供较均衡的路径形态诊断；通过稀疏触发将 CVF 额外评价控制在较低水平。本文据此将 CVF-AE 定位为强约束低空路径规划中兼顾可行性、轨迹形态诊断和评价开销的均衡型候选生成框架。
